%%\documentclass[pdflatex,sn-nature]{sn-jnl}% Style for submissions to Nature Portfolio journals
%%\documentclass[pdflatex,sn-basic]{sn-jnl}% Basic Springer Nature Reference Style/Chemistry Reference Style
\documentclass[pdflatex,sn-mathphys-num]{sn-jnl}% Math and Physical Sciences Numbered Reference Style

%%%% Standard Packages
%%<additional latex packages if required can be included here>
\usepackage{graphicx}%
\usepackage{multirow}%
\usepackage{amsmath,amssymb,amsfonts}%
\usepackage{amsthm}%
\usepackage{mathrsfs}%
\usepackage[title]{appendix}%
\usepackage{xcolor}%
\usepackage{textcomp}%
\usepackage{manyfoot}%
\usepackage{booktabs}%
\usepackage{tikz}
\usetikzlibrary{decorations}
\usetikzlibrary{decorations.pathreplacing}
\usepackage{algorithm}%
\usepackage{algorithmicx}%
\usepackage{algpseudocode}%
\usepackage{listings}%
\usepackage{xr-hyper}
\usepackage{lineno}
\usepackage{caption}
\externaldocument[S]{dunkelflaute_extended_data_figures}

%%%%

\begin{document}

\linenumbers


\title[Article Title]{Global risks of \textcolor{red}{wind-solar} energy droughts under climate change}

%%=============================================================%%
%% GivenName	-> \fnm{Joergen W.}
%% Particle	-> \spfx{van der} -> surname prefix
%% FamilyName	-> \sur{Ploeg}
%% Suffix	-> \sfx{IV}
%% \author*[1,2]{\fnm{Joergen W.} \spfx{van der} \sur{Ploeg} 
%%  \sfx{IV}}\email{iauthor@gmail.com}
%%=============================================================%%

\author*[1]{\fnm{Colin} \sur{Lenoble}}\email{colin.lenoble@universite-paris-saclay.fr}

\author[2,3]{\fnm{Mathias} \sur{Valla}}

\author[4]{\fnm{Jonas} \sur{Schmitt}}

\author[1]{\fnm{Céline} \sur{Guivarch}}

\author[1]{\fnm{Améline} \sur{Vallet}}

\author[5,6]{\fnm{Pauline} \sur{Rivoire}}

\author[7,8]{\fnm{Kevin} \sur{Schwarzwald}}

\author[9,7]{\fnm{Kai} \sur{Kornhuber}}

{\small 
\affil[1]{\orgdiv{AgroParisTech, Université Paris-Saclay, CIRAD, CNRS, EHESS, ENPC, Institut Polytechnique de Paris}, \orgname{CIRED}, \city{Nogent-sur-Marne}, \postcode{94130}, \country{France}}

\affil[2]{\orgdiv{I2M}, \orgname{Aix-Marseille Université}, \orgaddress{\city{Marseille}, \country{France}}}

\affil[3]{\orgdiv{Chaire DIALog}, \orgname{Institut Louis Bachelier}, \orgaddress{\city{Paris}, \country{France}}}

\affil[4]{\orgdiv{Agricultural Economics and Policy Group}, \orgname{ETH Zurich}, \orgaddress{ \city{Zurich}, \country{Switzerland}}}

\affil[5]{\orgname{University of Lausanne}, \orgaddress{ \city{Lausanne}, \country{Switzerland}}}

\affil[6]{\orgname{Ecole Polytechnique Fédérale de Lausanne}, \orgaddress{ \city{Sion}, \country{Switzerland}}}

\affil[7]{\orgdiv{Lamont-Doherty Earth Observatory}, \orgname{Columbia University}, \orgaddress{\city{Palisades, NY}, \country{United States}}}

\affil[8]{\orgname{Columbia Climate School}, \orgaddress{ \city{New York}, \country{United States}}}

\affil[9]{\orgname{International Institute for Applied Systems Analysis, Environment, Climate and Energy Program}, \orgaddress{ \city{Laxenburg}, \country{Austria}}}
}



\keywords{energy security, compound events, climate change mitigation, energy transition, climate change impact}



\abstract{
Solar and wind energy are central to the energy transition, but \textcolor{red}{solar and wind} energy droughts (\textcolor{red}{SWEDs}), \textit{i.e.} weather conditions associated with low wind and solar power, pose significant risks to power systems relying heavily on renewables. Future changes in \textcolor{red}{SWEDs} and their interaction with temperature extremes remain poorly understood at the global scale. Using up to 32 simulations from 14 CMIP6 models, we provide a global analysis of current and future \textcolor{red}{SWED} risks under four warming levels and a rigorous treatment of uncertainties. Globally, \textcolor{red}{SWED} severity has increased by 4.9\% over the last 40 years, and future projections follow an increasingly polarized spatial pattern: +8.5\% under 1.5°C warming, with milder average increases (+4.7\%) under 2°C as more regions experience decreasing \textcolor{red}{SWED} severity. Climate models struggle to reproduce observed trends in some regions, highlighting the need to better account for model uncertainty. Critically, regions most exposed to \textcolor{red}{SWEDs} also face temperature-driven demand shifts, challenging the resilience of renewable-oriented electricity systems. 
}

\maketitle
\clearpage
\section{Introduction}

A transition from fossil fuel-based energy systems towards renewable energy sources is essential for limiting greenhouse gas emissions and mitigating global warming \cite{IEA2024, ipcc2023synthesis} consistent with the Paris Agreement \cite{UNFCCC_ParisAgreement_2018}. Solar and wind energy have rapidly expanded in global electricity generation, outpacing installation rate projections due to technological advancements and decreased costs \cite{IEA2025GlobalEnergyReview, creutzigUnderestimatedPotentialSolar2017}.
However, solar and wind generation depend on weather conditions. The joint occurrence of low wind and cloudy weather conditions creates so-called \textcolor{red}{solar and wind} energy droughts (\textcolor{red}{SWEDs}) \cite{Rouges2025WeatherRegimes}. As solar and wind energy sources are expected to play a decisive role in low-carbon transition scenarios, periods of combined low production of wind and solar \textcolor{red}{and of net load represent crucial design parameters of renewable-dominated power systems.} Hence, a rigorous assessment of \textcolor{red}{SWEDs} and a transparent communication of uncertainties is particularly important for energy policy and planning \cite{Kornuberetal2025,wilczakWindSolarEnergy2025}. 

The existing literature has begun to address these challenges, but key gaps remain. First, while previous global studies have examined the impacts of climate change on solar and wind power separately, the combined effect of climate change on renewable energy availability remains poorly understood \cite{wangRisingWorldwideChallenges2025, liuClimateChangeImpacts2023}. Second, most existing studies focus on developed regions, especially North America and Europe \cite{Bracken2024, brownMeteorologyClimatologyHistorical2021, singhEnhancedSolarWind2024, RAYNAUD2018578, oteroCopulabasedAssessmentRenewable2022, kapicaPotentialImpactClimate2024}, limiting our understanding of climate change impacts at the global scale. These studies use heterogeneous climate datasets: some studies rely on historical reanalysis  \cite{RAYNAUD2018578}, others on CMIP6 simulations \cite{kapicaPotentialImpactClimate2024}, or large ensemble simulations \cite{vanderwielMeteorologicalConditionsLeading2019}. Third, future changes in wind speed and solar radiation are difficult to assess, because climate models do not fully capture how large-scale atmospheric circulation will respond to warming \cite{Shepherd2014}, and because persistent cloud-related biases affect the accuracy of regional solar radiation projections \cite{huangDiagnosingRadiationBiases2023}. \textcolor{red}{Because GCMs remain among the most useful tools available to anticipate climate change effects, evaluating their projections and providing an honest communication of uncertainty is essential for choosing no-regret adaptation strategies for energy grids \cite{douvilleAllClimateModels}.} \textcolor{red}{Indeed}, the effects of climate change on \textcolor{red}{SWEDs} are often quantified using only a few simulations \cite{brackenIntensifyingRenewableEnergy2025, wangRisingWorldwideChallenges2025}, and without comparing model estimates with observations (e.g., over a historical reference period)\textcolor{red}{, neglecting uncertainty in projected impacts} \cite{wilczakWindSolarEnergy2025}\textcolor{red}{. Including such evaluation and uncertainty quantification can inform both climatologists, by highlighting regions of high uncertainty in wind-solar energy droughts, and energy economists, by identifying regions where projections show robustness in the short term.} 


Periods of low wind and solar generation are particularly critical when they coincide with high electricity demand, giving rise to \textcolor{red}{solar and wind budget} droughts (\textcolor{red}{SWBDs}). Climate change affects both sides of this balance: it alters the frequency and severity of \textcolor{red}{SWEDs} while simultaneously shifting the demand for heating and cooling, with pronounced regional contrasts between low- and middle-latitude countries and northern regions \cite{lizanaGlobalGriddedDataset2026, wenzNorthSouthPolarization2017}. Although recent studies report an increasing severity of \textcolor{red}{SWBDs} in the past \cite{zhengClimateChangeImpacts2024}, \textcolor{red}{SWBD} projections under climate change scenarios are scarce. Existing projections of \textcolor{red}{SWBDs} tend to be either region-specific or centered on solely solar- or wind-based energy systems, often overlooking the interactions between wind and solar resources that can help mitigate variability and decrease generation intermittency \cite{juraszComplementarityResourceDroughts2021, wangRisingWorldwideChallenges2025, liuClimateChangeImpacts2023}.

Here we provide a global assessment of \textcolor{red}{SWED} and \textcolor{red}{SWBD} characteristics and their co-evolution under historical conditions and future global warming levels. Using globally gridded reanalysis  and multi-model climate simulations to account for uncertainties we address three research questions: (1) what are the spatial and temporal characteristics of \textcolor{red}{SWEDs} at the global scale? (2) How do \textcolor{red}{SWEDs} characteristics evolve under climate change? (3) To what extent do changes in \textcolor{red}{SWEDs} exacerbate \textcolor{red}{SWBDs}? 

We defined \textcolor{red}{SWEDs} as days of compound low wind and low solar generation below the 10th percentile of a reference period centered on a global warming of 0.61°C above pre-industrial levels (Extended Data Fig.~\ref{edf:red_ebd_schema}a). The analysis of historical \textcolor{red}{SWEDs} is based on \textcolor{red}{ERA5}, and projections use CMIP6 simulations \cite{cucchiWFDE5BiasadjustedERA52020, gwlcmip6, Bracken2024, oteroCopulabasedAssessmentRenewable2022}. We use up to 32 simulations from 14 CMIP6 models multivariately bias-aligned with our reanalysis \cite{cannonMultivariateQuantileMapping2018}. We characterize \textcolor{red}{SWEDs} by their frequency, duration, and intensity \textcolor{red}{\cite{wilczakWindSolarEnergy2025}}, and introduce an aggregated metric, annual \textcolor{red}{SWED} severity, that combines these three dimensions. Its interpretation is analogous to heating or cooling degree days, but applied to renewable energy capacity factors. The capacity factor is the ratio of actual electricity output to the maximum possible output at full load. Finally, we estimate the interaction between \textcolor{red}{SWEDs} and demand using a stylized temperature-based demand model, which provides a climate-conditioned indicator of \textcolor{red}{SWBD} changes in systems with 50\% renewable energy  \cite{liuClimateChangeImpacts2023}.

This article provides, to our knowledge, the most comprehensive assessment of past and future \textcolor{red}{SWEDs} and \textcolor{red}{SWBDs}, by explicitly accounting for model uncertainties with a multi-model framework. We propose an original approach to assess model performance by comparing simulated and observed trends in past annual \textcolor{red}{SWED} severity—an evaluation missing from most existing renewable climate-impact studies. Indeed, models fail to reproduce observed trends in Northeastern South America and Indonesia. Climate change both increases demand and decreases supply in low- and mid-latitude regions. In high-latitude regions, the decreases in supply are compensated due to reduced heating demand. These findings reveal a critical gap: the regions most exposed to compounding climate and energy risks remain the least studied, calling for targeted research and adaptation investment to support robust decarbonization scenarios.




\section{Results}\label{sec:results}


\subsection{Historical changes in \textcolor{red}{solar and wind} energy drought}\label{sec:historical}

Globally, annual \textcolor{red}{SWED} severity has increased by 4.9\% [95\% Confidence Interval (CI): -1.6\%, 13.8\%] between 1980-1999 and 2000-2019 (Fig. \ref{fig1:red_historical}a). Changes are spatially heterogeneous: four regions exhibit annual \textcolor{red}{SWED} severity changes exceeding 25\%: the Northern Amazon, Southern Africa, coastal regions in India and Southeast Asia (Fig.~\ref{fig1:red_historical}a.). In contrast, the Arabian Peninsula, Southern South America and Australia show decreasing annual \textcolor{red}{SWED} severity ranging from -10\% to -30\%.

\begin{figure}[hbt!]
\includegraphics[width=\textwidth]{resources/final_figs/fig1_valuebyalpha_slides_01.png}
\caption{\textbf{Historical changes in annual \textcolor{red}{SWED} severity.} 
(\textbf{a}) Relative change in annual \textcolor{red}{SWED} severity between 1980–1999 and 2000–2019 from \textcolor{red}{ERA5} reanalysis. Colour saturation indicates the baseline level of annual \textcolor{red}{SWED} severity. Grey areas denote regions with wind speeds below the cut-in threshold, which are excluded from the analysis. Low, medium, and high severity correspond to 0, 0.25, and 1 equivalent full load days, respectively (see Methods, Section~\ref{sec:method_REDind}).
(\textbf{b–f}) Regional time series of spatially averaged annual \textcolor{red}{SWED} severity for Western US (\textbf{b}), Northern Amazon (\textbf{c}), Egypt–Sudan (\textbf{d}), Kenya (\textbf{e}), and India (\textbf{f}). Annual \textcolor{red}{SWED} severity is expressed in equivalent full load days (blue crosses), with linear trends (red lines) and 95\% bootstrap confidence intervals (shaded). Boxes in panel (\textbf{a}) indicate the corresponding regions.}\label{fig1:red_historical}
\centering
\end{figure}

Furthermore, annual \textcolor{red}{SWED} severity is unevenly distributed across the world (Fig. \ref{fig1:red_historical}, Extended Data Fig.~\ref{edf:historical_climatology}d). We highlight several representative regions to illustrate different typologies of \textcolor{red}{SWED} variability and changes. Northern Amazon shows the strongest increasing trend worldwide while Egypt presents a decreasing trend.  Western U.S. and Egypt exhibit low interannual spread around the trend line. In contrast, Kenya exhibits high interannual variability (Fig. \ref{fig1:red_historical}b-e). This regional diversity reflects a broader pattern: interannual variability of annual \textcolor{red}{SWED} severity is unevenly distributed and higher in low- and middle-latitude regions (Extended Data Fig.~\ref{edf:historical_climatology}g). 

Historical solar capacity factor largely follows a latitudinal gradient, with the highest capacity factors in the subtropics, while wind capacity factors are distributed unevenly (Extended Data Fig.~\ref{edf:historical_climatology}e-f). Notably, forest areas have more frequent low-wind capacity factor days than other land uses due to low 10m wind at the canopy level, leading to higher annual severity (Fig. \ref{fig1:red_historical}a). 



About 5\% of the global land surface \textcolor{red}{shows} limited historical trend agreement with reanalysis under our validation criterion (i.e., fewer than half of the simulations yield trends statistically compatible with the observed annual \textcolor{red}{SWED} severity trend), \textcolor{red}{particularly} in Northern Amazon and Indonesia (Fig.~\ref{fig:model_skill}). \textcolor{red}{In contrast, regions such as India demonstrate robust agreement with observed trends, validating the model's capacity to capture regional variability in \textcolor{red}{SWED} dynamics where data constraints are less severe.} To focus subsequent analyses on regions where historical agreement is stronger under this criterion, these areas are masked in the main projection analyses. Applying a stricter performance criterion—requiring agreement in at least 25 of 32 simulations—excludes over 20\% of the land surface, mainly across Africa (Fig. \ref{fig:model_skill}a., yellow and red regions). 

 
\begin{figure}
  \centering
  % --- Top figure ---
  \includegraphics[width=\textwidth]{resources/final_figs/fig2_trend_validation.png}
    \caption{\textbf{Model skill in reproducing observed annual \textcolor{red}{SWED} severity trends.}
(\textbf{a}) Spatial agreement: share of the 32 simulations showing trends statistically consistent with \textcolor{red}{ERA5} reanalysis (1982–2019). 
(\textbf{b–c}) Bootstrap trend distributions for Northern Amazon and India respectively. Red: reanalysis; blue: individual simulations.}
    \label{fig:model_skill}
\end{figure}


\subsection{Changes in \textcolor{red}{solar and wind} energy droughts under global warming}\label{sec:projections}

Globally, multi-model mean annual \textcolor{red}{SWED} severity is projected to increase by 4.7\% [95\% CI: -2.2\%, +11.9\%] from 0.61°C to 2°C warming with significant regional differences (Fig.~\ref{gwl2_grid}a). Areas experiencing increasing annual \textcolor{red}{SWED} severity are partly differing with ones of the historical period. Several regions with weak or even decreasing historical trends projected on average by GCMs to face future increases, notably Central Asia, inland regions of the Indian subcontinent, the Arabian Peninsula, and parts of North America, as well as a northward shift of the increasing region over Africa. In parallel, more regions exhibit declining \textcolor{red}{SWED} severity under 2°C warming, particularly across Central America and the Mediterranean basin, in addition to Southern South America, which already shows historical decreases. 

\begin{figure}[hbt!]
\includegraphics[width=\textwidth]{resources/final_figs/fig3_projected_change_valuebyalpha_GWL2.png}
    \caption{\textbf{Projected changes in annual \textcolor{red}{SWED} severity at 2°C warming.}
(\textbf{a}) Multi-model ensemble mean relative change from the historical reference period (1982–2001, 0.61°C warming). Color saturation reflects reference annual \textcolor{red}{SWED} severity magnitude. Black areas indicate regions with limited historical trend agreement under the validation criterion and are therefore masked in the main analysis. Grey areas indicate regions with wind speed below the cut-in speed, which are excluded from the analysis.
(\textbf{b–e}) Regional annual \textcolor{red}{SWED} severity across ensemble members under four warming levels (0.61°C, 1.5°C, 2°C, 3°C) for Western US (\textbf{b}), South Africa (\textbf{c}), Kenya (\textbf{d}), and India (\textbf{e}). Dots: individual models; red line: weighted multi-model mean. Extended Data Figure 4 shows model agreement in trends. The global warming level 0.61 °C shows little variation across the simulations, as they have all been calibrated against the \textcolor{red}{ERA5} reference period (1982–2001).}
 \label{gwl2_grid}
\centering
\end{figure}

The response of annual \textcolor{red}{SWED} severity to warming is nonlinear. At 1.5°C, the global mean increase (+8.5\%, 95\% CI: +2.5\% to +14.8\%) exceeds the historical trend (+4.9\%), reflecting particularly large increases especially in India, Central Africa or North America (Extended Data Fig.~\ref{edf:annual_severity_gwl}). At 2°C, the global mean increase is lower (+4.7\%, 95\% CI: -2.2\% to +11.9\%), a difference that is not statistically distinguishable from the 1.5°C estimate. This reflects a shift in the large zone of low-change pixels (±10\%, covering ~45\% of the globe) from slightly positive to slightly negative, while the most affected regions continue to intensify with few regions with changing patterns such as Iran. (see Fig.~\ref{gwl2_grid}d). The increases observed at 1.5°C and 2°C are substantially higher at 3°C (e.g., Kenya), which leads to a higher global mean (+13.6\%, 95\% CI: +2.7\%, +24.3\%), with particularly strong amplification at both extremes of the distribution, suggesting growing spatial contrasts. A bootstrap across GCM runs yields consistent results, confirming that the non-monotonic pattern stems from both regional compensation and inter-model differences. Notably, two of the three models showing the smallest changes at 1.5°C and 2°C warming levels are absent from our 3°C ensemble (19 runs), suggesting that the 3°C estimates may be biased toward higher-sensitivity models.


Uncertainty of projected changes under 2°C warming is substantially higher in low- and middle-latitude regions than at higher latitudes (Extended Data Fig.~\ref{edf:annual_severity_agreement}). Several models project significant \textcolor{red}{SWED} increases across large parts of the globe, while regions where the multi-model mean projects decreases or low changes, such as Central America, Southeast Asia, Eastern Asia or Central Europe are associated with high inter-model uncertainty, both at 2°C and 3°C warming. Regions showing the strongest reversals between 1.5°C and 2°C, notably Iran, Western Africa, and South Australia, illustrate two distinct sources of uncertainty: high internal variability dominates in Iran and South Australia (Extended Data Fig.~\ref{edf:annual_severity_agreement}d), while inter-model disagreement is an important driver in Western Africa (Extended Data Fig.~\ref{edf:annual_severity_agreement}a–b). When aggregating grid-cell results to the regional scale, decreases in \textcolor{red}{SWED} severity are attenuated suggesting spatial compensations within the regions explaining part of the differences between Extended Data Figs.~\ref{edf:annual_severity_gwl} and \ref{edf:annual_severity_agreement}.

\subsection{Increasing \textcolor{red}{solar and wind budget} droughts under global warming}\label{sec:temporal_mismatch}


\begin{figure}[]
\includegraphics[width=\textwidth]{resources/final_figs/fig4_main_gwl_maps.png}
\caption{\textbf{Multi-model ensemble mean change in \textcolor{red}{solar and wind budget} droughts (\textcolor{red}{SWBDs}) under different global warming levels.} 
 Relative changes in \textcolor{red}{SWBDs} under (\textbf{a}) 2°C, (\textbf{b}) 1.5°C, and (\textbf{c}) 3°C global warming levels relative to the reference period (0.61°C). \textcolor{red}{SWBDs} are derived from a stylized residual load framework in which temperature-driven electricity demand is compared with renewable supply from existing wind and solar capacities scaled to meet 50\% of mean demand over the reference period. Events correspond to days when residual load exceeds the 99\textsuperscript{th} percentile of the reference-period distribution. Black regions are excluded based on trend-validation criteria, while hatched regions indicate areas without renewable capacities as of 2025. See Methods (Section~\ref{sec:methods}) for a detailed description of the model and event definition.}
\label{fig:supply_global_map}
\centering
\end{figure}

To assess \textcolor{red}{SWBD} changes across warming levels, we use a residual load approach linking temperature-driven demand with wind and solar supply (assuming 50\% renewable penetration in every region with current or planned renewable capacity). The effect of climate change on \textcolor{red}{SWBDs} varies regionally and differs in spatial pattern from changes in annual \textcolor{red}{SWED} severity (Fig.~\ref{fig:supply_global_map}). Of the 407 studied regions with existing renewable capacities, 189 experience an increase in multi-model mean \textcolor{red}{SWBDs} under 2°C warming, predominantly in low- and middle-latitude regions.  Conversely, \textcolor{red}{SWBDs} are projected to decrease in high-latitude regions.  

This latitudinal contrast strengthens with warming. Low- and middle-latitude regions face a progressive increase in \textcolor{red}{SWBDs} across warming levels, while high-latitude regions show minor reductions at 1.5°C and a substantial reduction at 3°C (Fig.~\ref{fig:supply_global_map}a–c). There is no major change in regional patterns across warming levels: regions exhibiting increasing (or decreasing) \textcolor{red}{SWBDs} at lower warming levels generally maintain the same trend at higher warming levels. Under 3°C warming, changes mainly reflect an amplification of existing trends, with only a few regions near transition zones shifting from improving to worsening conditions between 1.5°C and 3°C warming, such as Texas. 


To disentangle the drivers of these changes, we decompose \textcolor{red}{SWBDs} variations into demand- and supply-drivers (Fig.~\ref{fig:exp_regional_change} shows results for a representative subset of regions). Temperature-driven demand changes are, on average, four times larger than supply-driven changes and largely determine the direction of projected \textcolor{red}{SWBDs} shifts. In some regions such as Western China or Northern Europe supply changes also contribute substantially (Extended Data Fig.~\ref{edf:driver_decomposition}a). In high-latitude regions such as Ontario, Washington (USA), Sichuan, and Germany, declining heating demand under warming tends to reduce \textcolor{red}{SWBDs}. By contrast, in many low- and middle-latitude regions, including Egypt, the Ivory Coast, and Ecuador, rising cooling demand interacts with \textcolor{red}{SWEDs}, increasing \textcolor{red}{SWBDs}. In India and parts of Africa, increases in \textcolor{red}{SWBDs} across warming levels are driven by both demand and supply drivers, illustrating the compounding effect from demand and supply (Extended Data Fig.~\ref{edf:driver_decomposition}b-c). In contrast, although the Western U.S. is projected to experience increasing annual \textcolor{red}{SWED} severity, supply-driven changes do not translate into clear shifts in \textcolor{red}{SWBDs}. This reflects a seasonal decoupling: projected increases in \textcolor{red}{SWED} severity predominantly occur during winter, whereas \textcolor{red}{SWBDs} are more strongly driven by summer demand peaks for 50\% renewable energy systems. As a result, changes in renewable generation do not systematically coincide with periods of high demand, limiting their impact on \textcolor{red}{SWBDs}. 


\begin{figure}[hbt!]
\includegraphics[width=\textwidth]{resources/final_figs/fig5_main_dumbbell.png}
\caption{\textbf{Decomposition of the climate change effects on \textcolor{red}{solar and wind budget} droughts (\textcolor{red}{SWBDs})}. Multi-model mean of \textcolor{red}{SWBDs} relative change between 0.61°C and 2°C warming levels is represented in black and decomposed into two factors: (i) relative change in demand driver in purple, and (ii) relative change in renewable supply driver in yellow. The smaller dots represent the multi-model distributions for supply and demand effects. \textcolor{red}{SWBDs} are quantified as the cumulative exceedance of the daily residual load — difference between temperature-driven demand and current renewable supply - above its 99th percentile over the reference period. The x-axis is split at 800\%: the left panel uses a linear scale (-100\% to 800\%) and the right panel logarithmic scale (800\% to 1500\%). Regions were selected based on contrasting patterns between temperature-driven demand effects and renewable supply, as well as their geographical diversity.} \label{fig:exp_regional_change}
\centering
\end{figure}

Climate uncertainty propagates through both supply and demand drivers, shaping regional outcomes in contrasting ways (Fig.~\ref{fig:exp_regional_change}). Supply-driven uncertainty accounts for nearly half of the total uncertainty in \textcolor{red}{SWBDs} projections and even dominates in several regions, including South-East Asia, the Middle East, and North America (Extended Data Fig.~\ref{edf:driver_decomposition}d). Uncertainty is not distributed symmetrically across regions. It is substantially larger where rising cooling demand amplifies \textcolor{red}{SWBDs} than in regions where declining heating demand reduces their occurrence. 
 
Finally, as the total renewable share increases, \textcolor{red}{SWBD} events increasingly coincide with \textcolor{red}{SWED} days, suggesting a structural shift from demand-driven to supply-driven energy insecurity. This structural shift implies that the climate-driven increase in \textcolor{red}{SWBDs}, defined relative to system-specific thresholds, is reduced in high-renewable-penetration mixes, as the demand effect is partially offset by the growing role of supply variability (Extended Data Fig.~\ref{edf:re_share_sensitivity}). 


\clearpage


\section{Discussion}\label{sec:discussion}

To our knowledge, this study is the first of its kind to document changes in annual \textcolor{red}{SWED} severity under consideration of climate driven local energy demand (assessed as \textcolor{red}{SWBD} event), at the global scale under both historical conditions and four future global warming levels. In doing so we find strong regional heterogeneity in future energy security risks. Regions in the Global South face the strongest compounding pressures : rising cooling demand, increasing \textcolor{red}{SWED} severity, and strong climate uncertainty. Yet they are under-represented in renewable-climate impact studies \cite{gangopadhyayRoleWindsolarHybrid2022, juraszComplementarityResourceDroughts2021}, highlighting a critical research and policy gap.

A central contribution of this work is to demonstrate that the effects of climate change on wind–solar energy systems cannot be inferred from separate analyses of wind and solar resources. When considered independently, projected changes in wind and solar droughts exhibit spatially anti-correlated patterns, suggesting compensatory effects \textcolor{red}{between solar and wind} \cite{wangRisingWorldwideChallenges2025, liuClimateChangeImpacts2023, duttaFutureSolarEnergy2022, quProlongedWindDroughts2025a}. \textcolor{red}{However, accounting for their temporal dependence reveals hotspots of compounding droughts missed by marginal analysis} \cite{Zscheischler2020}. The compound analysis presented here tells a different story: while some regions experience compensatory effects that limit \textcolor{red}{SWED} increases (\textit{e.g.} Europe, Western China), other regions see their \textcolor{red}{SWED} risk increase under the compound approach (\textit{e.g.} Central Asia, Western U.S., India, Africa) (Fig~\ref{fig1:red_historical} and Fig~\ref{gwl2_grid}). These findings demonstrate the essential nature of compound analysis for hybrid renewable-oriented energy systems, where climate change simultaneously affects both production factors. Our analysis at different global warming levels reveals non-linear effects of warming on \textcolor{red}{SWEDs} and \textcolor{red}{SWBDs}. Indeed, warming increases spatial polarization: global \textcolor{red}{SWED} rises by +8.5\% at 1.5°C warming, but only by +4.7\% at 2°C, as more regions experience \textcolor{red}{SWED} decreases.

Using a comprehensive ensemble of climate models, we provide projections of \textcolor{red}{SWEDs} and a stylized \textcolor{red}{SWBD} metric while applying a robustness filter in regions where historical agreement with observed annual \textcolor{red}{SWED} severity trends is limited. This approach, inspired by studies about sea surface temperature anomalies \cite{ipcc2021physical} and temperature extremes \cite{kornhuberGlobalEmergenceRegional2024}; is for the first time applied to \textcolor{red}{SWEDs}. It indicates limited historical \textcolor{red}{SWED} trend agreement in Northern Amazon (consistent with Dutta et al. \cite{duttaFutureSolarEnergy2022}) and Indonesia (Fig~\ref{fig:model_skill}). This result calls for a greater consideration of model uncertainty in the analysis of \textcolor{red}{SWEDs}. We note that our model evaluation is purely statistical; future work could evaluate the physical processes in models driving \textcolor{red}{SWEDs} and their trends, which may also shed light on the causes of poor model performances we see here. Possible causes of these issues could be related to known issues in model simulations of sea surface temperature anomalies or aerosols \cite{schumacherExacerbatedSummerEuropean2024, bartokProjectedChangesSurface2017a}. 


\textcolor{red}{This study estimates wind power for both onshore and offshore turbines, extrapolating 10\,m wind speed to hub height (100\,m onshore, 120\,m offshore) using local wind shear coefficients fitted on historical 100\,m ERA5 wind speed, rather than a uniform global shear exponent, following standard practice in the literature given the coarse spatial resolution of GCMs \cite{quProlongedWindDroughts2025a}. This local calibration avoids the systematic underestimation of wind speed, and hence of capacity factors, particularly in coastal areas, that results from applying a uniform wind shear coefficient in climate-change impact assessments of wind-solar production.} For instance, our study shows an average wind capacity factor of 0.15 for the UK, compared to 0.25–0.4 reported by renewable energy companies \cite{DigestUKEnergy2025}. However, since onshore installations currently represent $90\%$ of total wind capacity globally ($80\%$ by 2050 \cite{IEA2024}), this simplification has a limited impact on our results.

Projecting the future spatial distribution of capacity factors was particularly challenging. Indeed, most regions lack spatially explicit scenarios of renewable energy deployment (as e-Highway 2050 for Europe \cite{sanchis2015}). In the absence of consistent global data on future deployment, particularly for developing countries, we used historical capacity factors as a proxy for resource quality and as aggregation weights. This approach contrasts with existing literature, in which deployment is often expected to occur in areas with no land use constraints (\textit{e.g}., \cite{wangRisingWorldwideChallenges2025, zhengStrategiesClimateresilientGlobal2025}). We compared our approach with simple surface area averages and we found that the difference is minor for solar energy but more pronounced for wind energy, particularly in Colombia and Chile (Extended Data Fig.~\ref{edf:aggregation_comparison}).  

The objective of this study was to isolate the effects of climate change on renewable-oriented energy systems, without addressing flexibility constraints associated with high renewable shares. We therefore assume a stylized system with a uniform renewable penetration of 50\% across regions, consistent with end-of-century SSP2 scenarios \cite{byersAR6ScenariosDatabase2022}, and focus on wind and solar energy only following Liu et al. \cite{liuClimateChangeImpacts2023}, without modeling how residual demand is met. Accordingly, the \textcolor{red}{SWBD} results should be interpreted as climate-conditioned stress indicators rather than full projections of future system adequacy. Sensitivity analyses across different penetration levels (25\%, 50\%, 75\%) and alternative wind-solar mixes, based on current and projected capacities, show limited impacts on the overall direction of \textcolor{red}{SWBD} changes (Extended Data Figs.~\ref{edf:re_share_sensitivity} and \ref{edf:mix_vs_warming}).



The choice of spatial and temporal scales strongly influences assessments of climate impacts on energy systems. This study was conducted at a regional (subnational) scale, which reveals intra country heterogeneity, in large countries such as the United States, China, and Russia, that is not captured in national-scale analyses \cite{liuClimateChangeImpacts2023, zhengStrategiesClimateresilientGlobal2025}.
We used daily data, consistent with multi-model ensemble studies \cite{oteroCopulabasedAssessmentRenewable2022, vanderwielMeteorologicalConditionsLeading2019}. While some studies rely on seasonal thresholds to detect \textcolor{red}{SWEDs}, these emphasize relative anomalies rather than absolute production deficits, which are more relevant for supply–demand balances. We therefore adopt an annual threshold, following common practice \cite{wilczakWindSolarEnergy2025}, and complement it with a sensitivity analysis. The sensitivity analysis confirms that our conclusions are robust to different capacity factor and residual load thresholds (Extended Data Fig.~\ref{edf:threshold_sensitivity}).


In most countries, changes in demand determine the overall direction of change in \textcolor{red}{SWBDs}. We chose a stylized demand model, with uniform renewable penetration and responses to heating and cooling to isolate the climate impact on both demand and supply for renewable oriented systems. This simplification is consistent with comparable studies \cite{liuClimateChangeImpacts2023, zhengStrategiesClimateresilientGlobal2025}. Sensitivity tests reveal low variations in the global direction of \textcolor{red}{SWBDs} changes (Extended Data Figs.~\ref{edf:re_share_sensitivity} and \ref{edf:demand_sensitivity}). \textcolor{red}{Our approach focuses exclusively on how changes in heating and cooling degree days affect electricity demand, deliberately excluding other key drivers such as demographic trends, economic development, and differential air-conditioning adoption curves \cite{vanruijvenAmplificationFutureEnergy2019b, davisAirConditioningGlobal2021}. This simplification carries important implications: by applying a uniform, synthetic demand response across regions, we do not capture population growth or the income-related gradient in thermosensitive response between low- and high-income countries [REF: Estimating a social cost of carbon for global energy consumption], likely leading us to overestimate SWBDs, especially at low global warming levels. However, air-conditioning adoption is accelerating in developing countries [REF: Accelerating room air conditioner efficiency in India: Grid, economic, and policy implications through 2035], and since high temperatures pose a health threat regardless of income level, a uniform approach serves as a conservative estimate that ensures cooling needs are not underestimated across all regions.}  


Our results support renewable deployment strategies already engaged in several regions, including the EU and Western China (Fig~\ref{gwl2_grid}a.), where \textcolor{red}{SWEDs} and \textcolor{red}{SWBDs} are projected to remain low or decrease. In contrast, regions in Africa and India may face growing renewable energy risks due to rising cooling demand and increasing \textcolor{red}{SWED} severity. Given the high levels of uncertainty in these regions, targeted region-specific studies are needed. Examining adaptation options such as storage, interconnections, and demand-side measures in high-risk regions is paramount, but requires finer temporal scales to capture intraday solar variations and accurately assess storage needs. Our multi-model framework, validated against historical data, provides a robust basis for prioritizing such analyses as demonstrated by the work from Zheng et al \cite{zhengStrategiesClimateresilientGlobal2025}. Finally, this work contributes to the growing literature on energy justice \cite{sovacoolEnergyJusticeConceptual2015, cherpConceptEnergySecurity2014}: regions facing the highest climate risks have contributed least to historical emissions, exhibit the highest energy poverty, and remain the least studied, highlighting the need to redirect research and investment toward these vulnerable regions.



\clearpage

\section{Methods and data}\label{sec:methods}

\subsection{Climate data}\label{sec:climate_data}


We used \textcolor{red}{the ERA5 reanalysis dataset, for the period 1982--2021,} to estimate historical solar and wind capacity factors \textcolor{red}{\cite{cds_era5_daily_2024}}.

For future projections, we used daily 2-meter maximum and mean temperature, surface downwelling solar (shortwave) radiation, and 10-meter zonal and meridional wind speeds (\textbf{tasmax}, \textbf{tas}, \textbf{rsds}, \textbf{uas}, and \textbf{vas}) from 14 GCMs from the 6th generation of the coupled model intercomparison project (CMIP6; \cite{eyring2016overview}). 

Following e.g. \cite{seneviratneRegionalClimateSensitivity2020, batibenizCountriesMostExposed2023}, we considered GCM output at particular Global Warming Levels (GWLs), defined as global mean surface temperature warming over a reference period. For each GCM ensemble member and GWL, we calculated a 20-year rolling average of  global mean surface temperature, and took the anomaly of that time series vs. each simulation's 1850-1900 average  global mean surface temperature. For each GWL, we found the first year that this smoothed time series crosses that GWL, and define the 20 years around it as the climate for that particular GWL. We used GWL 0.61°C (corresponding to 1982-2001 in the observed climate, as estimated using the HadCRUT5 global temperature product \cite{moriceUpdatedAssessmentNearSurface2021}) as a reference time period between models and observations, and GWLs 1.5°C, 2°C, and 3°C to examine future conditions. 

Using GWLs comes with two advantages over comparing GCMs using the calendar years of their simulations. Firstly, normalizing changes by GWL removes the dependence of climate changes on a particular forcing scenario. Forcing scenarios such as the Shared Socioeconomic Pathways (SSPs) are meant to be storylines rather than statistical samples of future socioeconomic conditions (e.g., \cite{riahiSharedSocioeconomicPathways2017a}); showing results by GWL thus allows a reader to interpret findings based on their own priors about timescales of earth system change. Secondly, normalizing changes by GWL reduces the influence of the so-called ``hot-model" problem, following arguments in which several GCMs from the CMIP6 generation may have unrealistically high climate sensitivities (i.e., \cite{hausfatherClimateSimulationsRecognize2022}). In both cases, showing results by GWL assumes that statistics of climate patterns depend on the magnitude of global warming (i.e,.  global mean surface temperature) but not the timing of when the earth system reaches that level of global warming; an assumption shared by the IPCC 6th Assessment Report \cite{IPCC2023SPM}.

Projections of future climate conditions at a particular GWL are affected by model uncertainty between GCMs (though differences between model climate sensitivities are suppressed when normalizing by GWLs) and the internal variability from the chaotic nature of the climate system \cite{hawkinsPotentialNarrowUncertainty2009}. For nonlinear transformations of climate variables in particular, internal variability is important to consider \cite{schwarzwaldImportanceInternalClimate2022a}. We characterized uncertainty across GCMs by using ensemble members from 14 GCMs. Given the constraints stemming from our use of computationally intensive multivariate bias correction techniques, we could not use all ensemble members from each GCM; we nevertheless used 10 members from CanESM5 and 9 members from MPI-ESM1-2-LR to estimate the magnitude of internal variability.
In total, we thus used 32 simulations from 14 GCMs for GWLs 0.61°C, 1.5°C, and 2°C, and 19 for GWL 3°C. GCM names, GWLs, and ensemble sizes are listed in Table~\ref{tab1}.
    
The period 1982–2001 served as the reference to apply a multivariate bias correction using the ``Multivariate Bias Correction with N-dimensional probability density function transform'' method \cite{cannonMultivariateQuantileMapping2018}. \textcolor{red}{The reanalysis was first aggregated to the respective spatial resolution of each GCM using the \texttt{conservative\_norm} function from \texttt{xesmf} before applying the bias correction.} This approach aligns not only the marginal distributions of the meteorological variables but also their copulas, preserving inter-variable dependencies relative to the observed reference period. However, it does not correct temporal properties such as trends or persistence: while the number of days above 30°C is matched, the duration of heatwaves (consecutive hot days) may differ between bias-corrected and observed datasets \cite{cannonMultivariateQuantileMapping2018}.

\begin{table}[h]
\caption{Number of members used, for each general circulation model (GCM, rows) and each global warming level (GWL, columns)}\label{tab1}%
\begin{tabular}{@{}lllll@{}}
\toprule
GCM & GWL 0.61°C  & GWL 1.5°C  & GWL 2°C & GWL 3°C\\
\midrule
ACCESS-CM2 & 1 & 1 & 1 & 1 \\
BCC-CSM2-MR & 1 & 1 & 1 & 0 \\
CanESM5 & 10 & 10 & 10 & 10 \\
CMCC-ESM2 & 1 & 1 & 1 & 1 \\
CNRM-CM6-1 & 1 & 1 & 1 & 1 \\
EC-Earth3-Veg-LR & 1 & 1 & 1 & 1 \\
GFDL-ESM4 & 1 & 1 & 1 & 0 \\
INM-CM5-0 & 1 & 1 & 1 & 0 \\
IPSL-CM6A-LR & 2 & 2 & 2 & 2 \\
KACE-1-0-G & 1 & 1 & 1 & 1 \\
MPI-ESM1-2-LR & 9 & 9 & 9 & 0 \\
MRI-ESM2-0 & 1 & 1 & 1 & 0 \\
NorESM2-MM & 1 & 1 & 1 & 0 \\
TaiESM1 & 1 & 1 & 1 & 1 \\
\botrule
\end{tabular}

\end{table}



\subsection{Administrative data}\label{adm_data}

To analyze \textcolor{red}{SWED} at the administrative level, we aggregated renewable energy capacity factors at the national scale. As not all regions have renewable capacities and to assess the effects of changes in the renewable energy mix, we decided to aggregate solar and wind distributions based on the geographic annual potentials as in \cite{huImplicationsParisproofScenario2023}\textcolor{red}{, using area-weighting from the \texttt{xagg} package}. For a subset of 13 large countries (\textit{Argentina, Australia, Brazil, Canada, China, Democratic Republic of the Congo, India, Indonesia, Kazakhstan, Mexico, Mongolia, Russia, United States}), we used administrative level 1 regions from the Natural Earth shapefile to ensure a more consistent regional size across countries. The corresponding shapefile is publicly available (open data; see repository).

\subsection{Renewable energy infrastructure data}\label{sec:infrastr_data}

Information on the location, capacity, and operational status of solar and wind farms were obtained from \cite{energyinfra}. Based on this dataset, we constructed two indicators of the solar-wind energy mix:
\begin{enumerate}
\item One based solely on currently operational renewable energy capacities (referred to as current mix).
\item Another that includes both existing capacities and those from announced or under construction projects (referred to as future mix).
\end{enumerate}



\subsection{Capacity factors estimation}\label{sec:cf_estimation}

Solar and wind capacity factors ($SCF$ and $WCF$) were computed following \cite{vanderwielMeteorologicalConditionsLeading2019} from meteorological variables only. We used daily time series at the grid-cell level and defined $SCF_{c,t}$ and $WCF_{c,t}$ as daily mean capacity factors (dimensionless, in $[0,1]$).\\
\\
\textcolor{red}{For onshore wind, daily mean hub-height wind speed (100\,m onshore) was estimated from 10\,m winds using a power-law profile \cite{Peterson1978}, with the shear exponent fitted locally using historical 10 and 100\,m ERA5 wind speeds. Wind  : }

\begin{equation}
    {\color{red}w_{100,c,t}=\left(\tfrac{100}{10}\right)^{\alpha_{c}} \times \sqrt{u_{10,c,t}^2+v_{10,c,t}^2}, \qquad \alpha_{c}\ \text{local wind shear coefficient}}
\end{equation}
where $u_{10}$ and $v_{10}$ were daily mean zonal and meridional wind components (m\,s$^{-1}$). The wind capacity factor was obtained from a three-regime turbine power curve,
\begin{equation}
    WCF_{c,t}=
\begin{cases}
0, & w<w_{\mathrm{ci}} \ \text{or}\ w\ge w_{\mathrm{co}},\\[3pt]
\dfrac{w^{3}-w_{\mathrm{ci}}^{3}}{w_{\mathrm{r}}^{3}-w_{\mathrm{ci}}^{3}}, & w_{\mathrm{ci}}\le w< w_{\mathrm{r}},\\[8pt]
1, & w_{\mathrm{r}}\le w< w_{\mathrm{co}},
\end{cases}
\qquad w=w_{80,c,t},
\end{equation}
with $(w_{\mathrm{ci}}, w_{\mathrm{r}}, w_{\mathrm{co}})=(3.5,13,25)$\,m\,s$^{-1}$, representing the cut-in, rated, and cut-out wind speeds, respectively \cite{vanderwielMeteorologicalConditionsLeading2019}.\\
\\
For solar PV, we used daily mean surface solar downward radiation $G_{c,t}$ (W\,m$^{-2}$) and a temperature-dependent performance ratio. The performance ratio was defined as
\begin{equation}
    PR_{c,t}=1+\gamma\left(T_{\mathrm{cell},c,t}-T_{\mathrm{ref}}\right),
    \qquad T_{\mathrm{ref}}=25^{\circ}\mathrm{C},\ \gamma=-0.005^\circ\mathrm{C}^{-1}, 
\end{equation}
with PV cell temperature (in $^{\circ}$C) estimated as
\begin{equation}
T_{\mathrm{cell},c,t}
= c_1+c_2\,\frac{T_{\max,c,t}+T_{c,t}}{2}+c_3\,G_{c,t}+c_4\,w_{c,t},
\end{equation}
where $T_{c,t}$ was daily mean near-surface air temperature ($^{\circ}$C), $T_{\max,c,t}$ was daily maximum near-surface air temperature ($^{\circ}$C), and $w_{c,t}$ was daily mean near-surface wind speed (m\,s$^{-1}$). $(c_1,c_2,c_3,c_4)=\left(
4.3~^{\circ}\mathrm{C},\,
0.943,\,
0.028~^{\circ}\mathrm{C}\,(\mathrm{W\,m^{-2}})^{-1},\,
-1.528~^{\circ}\mathrm{C}\,(\mathrm{m\,s^{-1}})^{-1}
\right)$ were taken from \cite{PhotovoltaicModuleThermal}. The daily mean solar capacity factor was then computed as
\begin{equation}
    SCF_{c,t}=PR_{c,t}\,\frac{G_{c,t}}{G_{\mathrm{stc}}},\qquad G_{\mathrm{stc}}=1000\,\mathrm{W\,m^{-2}}.
\end{equation}

Regional capacity factors used reference-period annual-mean cell capacity factors as siting weights:
\begin{equation}
    SCF_{i,t}=\frac{\sum_{c\in\mathcal{C}_i} \bar{SCF}_c\,SCF_{c,t}}{\sum_{c\in\mathcal{C}_i}\bar{SCF}_c},\qquad
    WCF_{i,t}=\frac{\sum_{c\in\mathcal{C}_i} \bar{WCF}_c\,WCF_{c,t}}{\sum_{c\in\mathcal{C}_i}\bar{WCF}_c},
\end{equation}
where $\mathcal{C}_i$ was the set of grid cells in region $i$, and $\bar{SCF}_c$, $\bar{WCF}_c$ were cell-level annual means over the reference period.\\
\\
We \textcolor{red}{evaluated} aggregated country-level daily series against two external datasets, using both Pearson and Spearman correlations computed country-by-country and summarized as distributions. Against the ENTSO-E country-level renewable production time series (2016--2019), we found Spearman correlations of $0.79$ (solar) and $0.80$ (wind) \cite{entsoe_powerstats}. Against the Renewables.ninja production time series (which uses MERRA-2), we found Spearman correlations of $0.92$ (solar) and $0.89$ (wind) \cite{pfenningerLongtermPatternsEuropean2016}. Higher agreement with Renewables.ninja was expected because both products were reanalysis-driven meteorology-to-power reconstructions, whereas ENTSO-E reflected additional system-level effects not represented here (time-varying installed capacity during 2016--2019, curtailment, outages, offshore/onshore composition, plant siting within countries, and grid constraints). Pearson correlations were systematically lower than Spearman correlations, consistent with biases in amplitude and scaling even when day-to-day ordering was well captured. We therefore emphasized Spearman correlation because our metrics and analysis were based on relative thresholds (quantiles), for which rank agreement was the most relevant validation criterion.


\subsection{\textcolor{red}{SWED} indicators}\label{sec:method_REDind}
\textcolor{red}{SWED} were identified using a coincidence-below-threshold method \cite{wilczakWindSolarEnergy2025, oteroCopulabasedAssessmentRenewable2022, Bracken2024}.
An energy drought day was defined as a day with compounding low wind and low solar capacity factors. Wind and solar capacity factor thresholds ($WCF_{thr}$ and $SCF_{thr}$) were defined as the 10th percentiles of the non-zero daily capacity factors during the reference period (global warming level 0.61$^{\circ}$C compared to the pre-industrial period, aligned with 1982–2001).

A \textcolor{red}{SWED} event was defined as a series of consecutive days of energy drought. 
We compared their main characteristics—frequency, duration, and intensity—across different levels of global warming.
Frequency was defined as the number of \textcolor{red}{SWED} events per year, while duration referred to the average number of days per event within a given year. 
Daily intensity, for grid cell $i$ and day $t$, was defined as the daily expected capacity factor power shortfall. More specifically, intensity represented daily deficits between \textcolor{red}{the defined 10th-percentile threshold values ($WCF_{\mathrm{thr}}$, $SCF_{\mathrm{thr}}$)} and \textcolor{red}{the actual capacity factor}:
\begin{equation}
\mathrm{Int}_{i,t} = \bigl(WCF_{\mathrm{thr},i}-WCF_{i,t}\bigr)_{+}
+ \bigl(SCF_{\mathrm{thr},i}-SCF_{i,t}\bigr)_{+}.
\end{equation}
Then, for grid cell $i$ and year $y$, intensity was defined as the mean daily intensity per event day:
\begin{equation}
\text{Intensity}_{i,y}=\frac{1}{D_{i,y}}\sum_{t\in \mathcal{D}_{i,y}}\mathrm{Int}_{i,t},
\end{equation}
where $\mathcal{D}_{i,y}$ were \textcolor{red}{SWED} days in cell $i$ and year $y$, and $D_{i,y}=|\mathcal{D}_{i,y}|$ was the total number of \textcolor{red}{SWED} days.\\
\\
\noindent Drawing on the concept of cooling and heating degree days \cite{mirandaChangeCoolingDegree2023}, we developed an annual composite risk score, referred to as annual severity, that accounted for these three indicators:
\begin{equation}
\text{Annual Severity}_{i,y}
= \sum_{t \in \mathcal{D}_{i,y}} \mathrm{Int}_{i,t}
\end{equation}

Equivalently, using the definitions of frequency, duration, and intensity, this can be written as:

\begin{equation}
\text{Annual Severity}_i=\text{Frequency}_i\times\text{Duration}_i\times\text{Intensity}_i.
\end{equation}

For grid cell $i$, $\text{Annual Severity}_i$ represented the annual sum of renewable energy deficits during \textcolor{red}{SWED} events. Annual severity was expressed in equivalent full load days. 

Regions poleward of 68°N and 58°S were excluded due to artifacts in the duration metric, which became unstable under the near-permanent solar drought conditions characteristic of polar regions.

When computing the multi-model mean, we weighted simulations so that each GCM contributed equally, regardless of the number of ensemble members it provided. This prevented models with larger ensembles from disproportionately influencing the aggregated signal and ensured that trends reflected inter-model diversity rather than ensemble size. Moreover, as GCMs did not have the same resolution, we downscaled using a nearest-neighbor method applied to the capacity factors to map them to the reanalysis grid ($0.1^{\circ}\times 0.1^{\circ}$).

To quantify the global mean change in annual \textcolor{red}{SWED} severity between periods (or warming levels), we computed area-weighted spatial means using cosine-latitude weights, aggregating across all land pixels passing the trend validation criteria. The relative change was derived by dividing the difference between projected and baseline weighted multi-model means by the baseline weighted mean. Uncertainty bounds were estimated using a spatial block bootstrap procedure ($n = 1{,}000$ iterations) applied to the absolute change field. To account for spatial autocorrelation in the climate fields, bootstrap resampling was performed on contiguous $10^{\circ} \times 10^{\circ}$ blocks rather than individual pixels. The 95\% confidence intervals were derived from the 2.5th and 97.5th percentiles of the bootstrap distribution, scaled by the baseline global mean.


\subsection{Trend evaluation}\label{meth:trend}

We evaluated the ability of bias-corrected simulations to reproduce observed trends in annual \textcolor{red}{SWED} severity over the period 1982–2019. For each simulation, we selected a 38-year window whose first 20 years were centered on a GWL of 0.61°C (the period used to bias-correct projections). Observed trends were calculated from reaggregated reanalysis data at each grid cell and were compared with the corresponding trends from the 32 simulations over their aligned 38-year windows.

To ensure consistency in spatial scale, reanalysis data were first reaggregated to the native spatial resolution of each GCM before trend estimation. Uncertainty ranges for trend estimates in both observations and simulations were derived using a stationary bootstrap with 1,000 resamples and a block length of 3 years \citep{Wilks2006, politisStationaryBootstrap1994} to account for temporal autocorrelation in the annual severity time series. Trend compatibility was declared when the 95\% bootstrap confidence intervals overlapped between observations and simulations.

\subsection{Statistical significance} \label{meth:signif}

We assessed the statistical significance of changes by testing whether the distribution of our indicator differed between the historical climate ($\text{GWL } 0.61^\circ C$) and the $\text{GWL }2^\circ C$ using a non-parametric permutation test (see \cite{Fisher1935, Welch1990, Edgington2007}). Under the null hypothesis that realizations from both climates were exchangeable, we pooled all realizations, computed a test statistic \(T\) (difference in ensemble means; conclusions were robust to using medians), and approximated the null distribution by repeatedly reassigning climate labels at random and recomputing \(T\). The two-sided \(p\)-value was defined as the fraction of permutations with \(|T^\ast| \ge |T_{\text{obs}}|\). We preferred this non-parametric approach because it avoided strong distributional assumptions, yielded valid finite-sample inference under exchangeability, and was robust to skewness, outliers, and heteroskedasticity common in climate indicators; it also accommodated paired-member designs when applicable. We used $10000$ permutations with a fixed seed.

For computational efficiency, gridded maps were aggregated to the regional or country level, as described in Section~\ref{adm_data}. For multiple comparisons across regions or indicators, we controlled the false discovery rate (FDR) using the Benjamini--Hochberg procedure (\cite{Benjamini1995}) and declared significance at the resulting data-driven threshold, following recommendations from \cite{Wilks2016} and established practice for climate ensembles \cite{Zeman2022}. FDR control was applied per indicator across regions with a 0.2 threshold within each comparison.

For each region \(i\), global climate model (GCM) \(g\), and run \(r\), we encoded the outcome of the permutation tests as
\begin{equation}
    S_{g,r,i} =
    \begin{cases}
    +1, & \text{significant increase},\\
    -1, & \text{significant decrease},\\
    0,  & \text{otherwise},
    \end{cases}
\end{equation}
at the FDR-controlled significance level. We averaged runs within each GCM so that each model contributed equally regardless of its number of runs, $R_g$:
\begin{equation}
M_{gi} \;=\; \frac{1}{R_g}\sum_{r=1}^{R_g} S_{g,r,i},
\end{equation}
where $M_{gi}$ was the significance score of model $g$ in region $i$, averaged across its runs. We then averaged across models:
\begin{equation}
A_i \;=\; \frac{1}{G}\sum_{g=1}^{G} M_{g,i}
\;=\; \sum_{g=1}^{G}\sum_{r=1}^{R_g} \frac{1}{G\,R_g}\, S_{g,r,i}.
\end{equation}
The agreement index \(A_i\in[-1,1]\) summarized consensus on the direction of change: values near \(+1\) indicated widespread significant increases across models, values near \(-1\) indicated widespread significant decreases, and values near \(0\) indicated low or mixed agreement (see Extended Data Fig.~4).

\subsection{Uncertainty decomposition}\label{method:uncertainty}

We characterized the uncertainty in projections of \textcolor{red}{SWED} risk as the spread in outcomes across our ensemble of GCM runs. This uncertainty could be attributed to two sources: model uncertainty, arising from differences between GCMs due to modeling choices and resolution, and irreducible internal variability, stemming from the chaotic nature of the climate system \cite{schwarzwaldImportanceInternalClimate2022a, deserUncertaintyClimateChange2012}. Working with warming levels removed the scenario uncertainty that is part of the uncertainty triptych in climate science \cite{hawkinsPotentialNarrowUncertainty2009}. To understand the drivers of uncertainty in our projections, we estimated the relative contribution of each of these sources to the total variance. We followed \cite{lehnerPartitioningClimateProjection2020} and estimated uncertainties as follows.

We estimated model uncertainty as the variance across model means:
\begin{equation}
    M = Var_m\left(\frac{1}{N_{\mathrm{run}}}\sum_{r=1}^{N_{\mathrm{run}}} X_{m,r}\right)
\end{equation}
for a variable $X$ projected by run $r$ from model $m$. In practice, the models and run counts used were detailed in Table~\ref{tab1}.

We estimated internal variability as the mean variance across GCM runs:
\begin{equation}
    I = \frac{1}{N_{\mathrm{GCM}}}\sum_{m=1}^{N_{\mathrm{GCM}}}\left(Var_r\left(X_{m,r}\right)\right).
\end{equation}
Because of the computational cost of producing global multivariate bias-corrected and downscaled projections, internal variability was estimated using only two GCMs: CanESM5 and MPI-ESM1-2-LR, with 10 and 9 analyzed runs, respectively. In particular, $I + M$ was not guaranteed to equal the total variance when the number of runs per model differed.

\subsection{Energy balance droughts}\label{sec:meth:supply-demand}

Renewable capacities were extracted at the regional level (see Section~\ref{adm_data}). 
We constructed a time-varying renewable capacity factor by aggregating solar and wind capacity factors using the solar–wind mix ratio ($\alpha$, Section~\ref{sec:infrastr_data}):

\begin{equation}
    E_{i,t} = (\alpha  \, SCF_{i,t} + (1-\alpha)\, WCF_{i,t})\, RE_{\text{share}}
\end{equation}

Here, $E_{i,t}$ was a capacity-factor-based proxy for wind–solar generation availability (dimensionless). In practice, $E_{i,t}$ could be interpreted as wind–solar generation normalized by installed wind–solar capacity (up to a proportionality constant $RE_{\text{share}}$ set by the scaling described below). Importantly, $E_{i,t}$ did not represent total electricity supply, but only the wind–solar component.

The thermo-sensitive demand was derived from current population-weighted temperature time series as in \cite{zhengStrategiesClimateresilientGlobal2025}: 

\begin{equation}
    \text{D}_{i,t} = \hat D_i \, \left(1 + \delta_{HDD}\, (T_{\text{min}}-T_{i,t})_+ + \delta_{CDD}\, (T_{i,t} -T_{\text{max}})_+\right).
\end{equation}

where $\delta_{HDD}$ and $\delta_{CDD}$ were the heating and cooling demand responses, respectively, fixed as in \cite{zhengStrategiesClimateresilientGlobal2025} at $2.6\%\ ^\circ C^{-1}$ and $3.5\%\ ^\circ C^{-1}$, in agreement with developed-country regions ($4\%$ in the U.S. for cooling \cite{EIA2021AC}), assuming a uniform equipment rate of electric heating and air conditioning across regions. $T_{\text{min}}$ and $T_{\text{max}}$ were fixed globally at 12.5°C and 19.6°C, respectively. $\hat D_i$ corresponded to the baseline demand, and its calculation is explained below. This stylized formulation was intended to capture first-order, temperature-driven shifts in demand. Applying uniform response coefficients globally was a simplifying assumption, as heating and cooling elasticities vary substantially across regions with different climates, income levels, and equipment rates. 

We imposed long-term equilibrium over the reference period by requiring that renewable generation supplied a fixed average share, denoted $RE_{share}$, of temperature-adjusted demand:

\begin{equation}
    \mathbb{E}_t[E_{i,t}] = RE_{share} \, \mathbb{E}_t\![D_{i,t}]
\end{equation}

Substituting the temperature-based demand formulation yielded a closed-form expression for baseline demand $\hat D_i$, ensuring that wind and solar provided 50\% (see Extended Data Fig.~\ref{edf:re_share_sensitivity} for 25\% and 75\%) of average annual electricity demand over the reference period. This calibration defined a consistent renewable-oriented reference system in which wind–solar generation provided a fixed mean energy share $RE_{share}$, without specifying how the remaining fraction $(1-RE_{share})$ was met. The resulting demand profile was therefore stylized and not intended to represent observed national electricity systems; rather, it was designed to isolate the climatic contribution to supply–demand variability while preserving the temporal structure of renewable generation.

We therefore interpreted residual load as a proxy for the balancing or firm-capacity requirement induced by wind–solar variability rather than as a full system-dispatch or adequacy simulation.

Residual load represented the electricity demand remaining after accounting for renewable generation:

\begin{equation}
    RL_{i,t} = D_{i,t}-E_{i,t}
\end{equation} \label{eq:RL}

Similarly to annual severity, we quantified energy insecurity as the cumulative intensity of residual-load exceedances above the $99^{\mathrm{th}}$ percentile of the reference period (other thresholds are shown in Extended Data Fig.~10). We defined Energy Balance Droughts (\textcolor{red}{SWBDs}) as series of consecutive days during which the residual load exceeded its 99th percentile. This threshold approximated the joint exceedance probability of the double 10\% criterion used for \textcolor{red}{SWEDs}, assuming independence between wind and solar capacity factors. Relaxing this independence assumption would yield a region-specific threshold reflecting the local correlation structure between wind and solar resources. In this interpretation, \textcolor{red}{SWBDs} captured the frequency and magnitude of high residual-load stress days (i.e.\ days requiring unusually large non-wind/solar balancing), rather than realized shortages, which would additionally depend on the operational flexibility and adequacy of non-wind/solar resources.

To distinguish the contributions of demand and supply drivers to changes in \textcolor{red}{SWBDs}, we adopted a counterfactual decomposition approach based on empirical quantile mapping. For each simulation, we calculated the cumulative intensity of residual-load exceedances according to four scenarios: (A) the reference climate (warming level 0.61°C) for demand and supply, (B) the future warming level for both, (C) the future temperature distribution mapped onto the reference renewable generation, and (D) the reference temperature distribution mapped onto the future renewable generation. Scenarios (C) and (D) were constructed using empirical quantile mapping to align the source distribution with the target, ensuring that only the change in distribution of one factor was introduced while the other was kept at its baseline state. This approach preserved the temporal alignment between supply and demand, which would otherwise have been lost if they were derived from different warming levels. The demand driver was then isolated as the difference between scenario (C) and the reference (A), reflecting changes in residual load induced solely by variations in temperature distribution, while the renewable supply driver was isolated as the difference between scenario (D) and the reference. The sum of the two marginal effects approximately corresponded to the total change (scenario (B) minus (A)), with any residual capturing nonlinear interactions between the two factors.


\clearpage



\bmhead{Acknowledgements}
The authors would like to thank Dánnell Quesada for advising on the choice of reanalysis and bias correction methods.
This work emerged from the 3rd Como Training School on Compound Events Statistics and Physics of Compound Events\\
C.L. thanks the French Agence de l’Environnement et de la Maîtrise de l’Énergie (ADEME). M.V. acknowledges support from the Research Chair DIALog under the aegis of the Risk Foundation, a joint initiative of Université Paris-Dauphine PSL and CNP Assurances.

\section*{Data availability}

Reanalysis climate data used in this study are publicly available on \textcolor{red}{Copernicus Climate Data Store}.

The CMIP6 climate projections used in this study were accessed from the Pangeo cloud data archive. Global warming level (GWL) periods were identified using the warming levels as in Seneviratne \textit{et al} \cite{seneviratneRegionalClimateSensitivity2020}  and 20-year increments of data were then extracted and preprocessed accordingly. 

Information on the location, installed capacity, and operational status of solar and wind energy facilities was sourced from the Global Energy Monitor, a publicly accessible database available at \href{https://globalenergymonitor.org/}{globalenergymonitor.org}.

Administrative boundary data were obtained from the Global Administrative Areas database (GADM), available at \href{https://gadm.org/}{gadm.org}.
 

\section*{Code availability}

The code used to generate the results in this study will be made publicly available on GitLab upon publication. The repository will contain the scripts and instructions necessary to reproduce the analyses presented in this work. 
The analysis have been done using Python 3.10. The code supporting the analyses and some outputs are accessible in the Recherche Data Gouv temporary repository \href{https://entrepot.recherche.data.gouv.fr/privateurl.xhtml?token=0d562a4d-06f6-43d6-a0b5-eaa1275c1104
}{https://entrepot.recherche.data.gouv.fr/privateurl.xhtml?token=0d562a4d-06f6-43d6-a0b5-eaa1275c1104} and will be available on Gitlab, under the CC-NC-BY-SA license.


\section*{Author contributions}

K.K., P.R. and K.S. developed the original idea and contributed to the broader framing of the study. C.L. conceived the present study, designed the methodology and led the project. C.L., K.S. and M.V. curated the data. C.L. developed the analytical framework and code, carried out the main analyses, and generated the figures. M.V. contributed to code development, supporting analyses, project coordination and manuscript revision. All authors contributed to the interpretation of the results. C.L. wrote the first draft of the manuscript. C.G., C.L., M.V. and A.V. substantially revised the manuscript. All authors reviewed, edited and approved the final version of the manuscript.

\clearpage
\bibliography{sn-bibliography}% common bib file
%% if required, the content of .bbl file can be included here once bbl is generated
%%\input sn-article.bbl

\clearpage
\
\setcounter{figure}{0}
\renewcommand{\thefigure}{\arabic{figure}}
\graphicspath{{resources/extended_data_figures/}}


\captionsetup[figure]{name={Extended Data Fig.}, labelsep=period}

\section*{Extended Data Figures}


\vspace*{\fill}

\begin{figure}[H]
  \centering
  \includegraphics[width=\textwidth,keepaspectratio]{suppfig1_combined_red_ebd_schema.png}
  \caption{\textbf{Illustrative construction of \textcolor{red}{solar and wind} energy droughts (\textcolor{red}{SWEDs}) and \textcolor{red}{solar and wind budget} droughts (\textcolor{red}{SWBDs}) for Pakistan.}\\
  \textbf{a,} Example of a \textcolor{red}{SWED} diagnostic based on daily wind and solar capacity factors in 1994. Dashed lines mark the 10th-percentile thresholds used to define low-wind and low-solar conditions; shaded intervals indicate periods when both variables are simultaneously below threshold. The lower panel shows the daily compound deficit, computed as the sum of wind and solar shortfalls relative to threshold on \textcolor{red}{SWED} days.
  \textbf{b,} Example of an \textcolor{red}{SWBD} diagnostic for the same region and year. The upper panel shows normalised temperature-sensitive demand and renewable supply, the middle panel shows residual load (demand minus renewable supply) together with the \textcolor{red}{SWBD} threshold (dashed line; 95th percentile in the figure), and the lower panel shows daily residual-load exceedances. Shaded intervals indicate an \textcolor{red}{SWBD} event window.}
  \label{edf:red_ebd_schema}
\end{figure}

\vspace*{\fill}
\clearpage


% --------------------------------------------------
% Extended Data Fig. 2
% --------------------------------------------------
\vspace*{\fill}

\begin{figure}[H]
  \centering
  \includegraphics[width=\textwidth,keepaspectratio]{suppfig2_mean_variables_6panel.png}
  \caption{\textbf{Historical climatology of \textcolor{red}{SWED} characteristics and underlying renewable resources from reanalysis.}
  Spatial patterns of mean \textbf{a,} compound-event frequency, \textbf{b,} compound-event duration, \textbf{c,} compound-event intensity, \textbf{d,} annual \textcolor{red}{SWED} severity, \textbf{e,} wind capacity factor, \textbf{f,} solar capacity factor, and \textbf{g,} interannual variability (standard deviation) of annual \textcolor{red}{SWED} severity, estimated from \textcolor{red}{ERA5} reanalysis over the historical period. Gridpoint in grey denote areas excluded from the analysis due to wind speed below the cut-in speed.}
  \label{edf:historical_climatology}
\end{figure}

\vspace*{\fill}

% \section{Changes in \textcolor{red}{solar and wind} energy droughts under global warming}\label{supp:sec:projections}
% --------------------------------------------------
% Extended Data Fig. 3
% --------------------------------------------------

\vspace*{\fill}

\begin{figure}[H]
  \centering
  \includegraphics[width=\textwidth,keepaspectratio]{suppfig3_valuebyalpha_all_gwl.png}
  \caption{\textbf{Projected changes in annual \textcolor{red}{SWED} severity under three global warming levels.}
  Multi-model mean relative changes in annual \textcolor{red}{SWED} severity at \textbf{a,} 1.5$^\circ$C, \textbf{b,} 2.0$^\circ$C and \textbf{c,} 3.0$^\circ$C warming, each relative to the 0.61$^\circ$C reference warming level. Hue indicates the sign and magnitude of the projected change, whereas colour saturation indicates the reference annual-severity class (low, medium or high). Black regions are excluded by the trend-validation filter used in the main analysis. Dark grey grid points denote areas excluded from the analysis due to wind speed below the cut-in speed.}
  \label{edf:annual_severity_gwl}
\end{figure}

\vspace*{\fill}


% --------------------------------------------------
% Extended Data Fig. 4
% --------------------------------------------------

\vspace*{\fill}

\begin{figure}[H]
  \centering
  \includegraphics[width=\textwidth,keepaspectratio]{suppfig4_agreement_variability_combined.png}
  \caption{\textbf{Model agreement on changes in annual \textcolor{red}{SWED} severity and sources of projection uncertainty.} 
\textbf{(a, b)} For each region, colours indicate the fractions of weighted simulations showing a significant increase and decrease in annual \textcolor{red}{SWED} severity relative to the $0.61^\circ$C reference warming level (90\% confidence), under $2^\circ$C (\textbf{a}) and $3^\circ$C (\textbf{b}) warming. Red hues denote a majority of simulations projecting increases, blue hues decreases, and yellow hues mixed signals. Grey regions indicate weak or no agreement on significant change. Black regions are excluded by the trend-validation filter. 
\textbf{(c)} Total projection spread under $2^\circ$C warming (log scale), combining internal variability and model uncertainty. 
\textbf{(d)} Fraction of internal variability in total uncertainty ($I/(I+M)$) under $2^\circ$C warming: green hues indicate dominance of internal variability, and pink hues dominance of inter-model differences.}
  \label{edf:annual_severity_agreement}
\end{figure}

\vspace*{\fill}

\clearpage

% \section{Increasing \textcolor{red}{solar and wind budget} droughts under global warming}\label{supp:sec:temporal_mismatch}
% --------------------------------------------------
% Extended Data Fig. 5
% --------------------------------------------------

\vspace*{\fill}

\begin{figure}[H]
  \centering
  \includegraphics[width=\textwidth,keepaspectratio]{suppfig5_combined_driver_effects.png}
  \caption{\textbf{Decomposition of \textcolor{red}{SWBD} changes at 2.0$^\circ$C warming into renewable-supply and temperature-driven demand contributions.}
  \textbf{a,} Share of the absolute total effect attributable to renewable supply, computed as $|RE|/(|RE|+|Demand|)$. Values near 1 indicate predominantly supply-driven changes, whereas values near 0 indicate predominantly demand-driven changes.
  \textbf{b,} Multi-model mean renewable-supply effect on \textcolor{red}{SWBDs}. \textbf{c,} Multi-model mean demand effect on \textcolor{red}{SWBDs}. Positive values indicate larger \textcolor{red}{SWBDs}, and negative values smaller \textcolor{red}{SWBDs}, relative to the 0.61$^\circ$C reference.
  \textbf{d,} Share of the inter-model spread attributable to renewable supply, computed as $\sigma_{RE}/(\sigma_{RE}+\sigma_{Demand})$. Values near 1 indicate uncertainty dominated by renewable-supply changes.
  Black regions are excluded by the trend-validation filter; hashed regions indicate areas without current renewable capacities.}
  \label{edf:driver_decomposition}
\end{figure}

\vspace*{\fill}
\clearpage

% --------------------------------------------------
% Extended Data Fig. 6
% --------------------------------------------------


\vspace*{\fill}

\begin{figure}[H]
  \centering
  \includegraphics[width=\textwidth,keepaspectratio]{suppfig6_re_share_effect.png}
  \caption{\textbf{Sensitivity of \textcolor{red}{SWBD} changes to renewable penetration level.}
  Multi-model mean combined change in \textcolor{red}{SWBDs} for renewable penetration levels of 25\%, 50\% and 75\% of mean demand (rows) under global warming levels of 1.5$^\circ$C, 2.0$^\circ$C and 3.0$^\circ$C (columns). The regional wind--solar mix is held at the current mix used in the study, and \textcolor{red}{SWBDs} are defined from residual-load exceedances above the 99th percentile of the reference period. Warmer colours denote larger increases in \textcolor{red}{SWBDs} and greener colours decreases. Black regions are excluded by the trend-validation filter; hashed regions indicate missing data or zero current renewable capacity.}
  \label{edf:re_share_sensitivity}
\end{figure}

\vspace*{\fill}
\clearpage

% \section{Discussion}\label{supp:sec:discussion}
% --------------------------------------------------
% Extended Data Fig. 7
% --------------------------------------------------

\vspace*{\fill}

\begin{figure}[!htbp]
  \centering
  \includegraphics[width=\textwidth,keepaspectratio]{suppfig7_aggregation_comparison.png}
  \caption{\textbf{Comparison of regional capacity-factor aggregation methods.}
  Maps show the absolute difference in time-averaged solar (a) and wind (b) capacity factors between the potential-weighted mean (v1) and the area-weighted mean (v2). Scatter plots (c. and d.) compare v1 and v2 values across regions, with point colour indicating the absolute difference between methods; the dashed line shows the 1:1 relationship.}
  \label{edf:aggregation_comparison}
\end{figure}
\vspace*{\fill}
\clearpage


% --------------------------------------------------
% Extended Data Fig. 8
% --------------------------------------------------

\vspace*{\fill}

\begin{figure}[!htbp]
  \centering
  \includegraphics[width=\textwidth,keepaspectratio]{suppfig8_mix_gwl_effects.png}
  \caption{\textbf{Comparison of global-warming and renewable-mix effects on \textcolor{red}{SWBDs}.}
  % \textbf{a,} Scatter plot of the global-warming effect on \textcolor{red}{SWBDs} at 2.0$^\circ$C relative to the 0.61$^\circ$C reference level versus the renewable-mix effect associated with changing from the current to the future regional wind--solar mix, evaluated at the reference level. Blue markers indicate regions where the future mix has a larger wind share, and orange markers regions where it has a larger solar share. Axes are shown on symmetric logarithmic scales.
  \textbf{b,} Spatial classification of the combined effects. Dark red regions show increasing \textcolor{red}{SWBDs} from both warming and mix change; dark blue regions decreasing \textcolor{red}{SWBDs} from both; yellow regions increases from warming but decreases from mix change; and light blue regions decreases from warming but increases from mix change. Black regions are excluded by the trend-validation filter; hashed regions indicate areas without current renewable capacities.
  }
  \label{edf:mix_vs_warming}
\end{figure}

\vspace*{\fill}

\clearpage

% --------------------------------------------------
% Extended Data Fig. 9
% --------------------------------------------------

\vspace*{\fill}

\begin{figure}[H]
  \centering
  \includegraphics[width=\textwidth,keepaspectratio]{suppfig10_combined_threshold_sensitivity.png}
  \caption{\textbf{Sensitivity of \textcolor{red}{SWED} and \textcolor{red}{SWBD} changes to event-threshold choices.}
  The top row shows historical changes in annual \textcolor{red}{SWED} severity between 1980--1999 and 2000--2019 for \textcolor{red}{SWED} definitions based on joint wind--solar thresholds of \textbf{a,} 5th percentile, \textbf{b,} 10th percentile (reference) and \textbf{c,} 20th percentile. The bottom row shows the combined \textcolor{red}{SWBD} change at 2.0$^\circ$C warming for residual-load thresholds of \textbf{d,} 95th percentile, \textbf{e,} 99th percentile (reference) and \textbf{f,} 99.5th percentile. Red outlines mark the reference settings used in the main analysis. Masked regions follow the exclusions applied in the corresponding \textcolor{red}{SWED} and \textcolor{red}{SWBD} calculations.}
  \label{edf:threshold_sensitivity}
\end{figure}

\clearpage

% \section{Methods and Data}\label{supp:sec:methods}
% --------------------------------------------------
% Extended Data Fig. 10
% --------------------------------------------------

\vspace*{\fill}

\begin{figure}[!htbp]
  \centering
  \includegraphics[width=\textwidth,keepaspectratio]{suppfig9_demand_sensitivity.png}
  \caption{\textbf{Sensitivity of \textcolor{red}{SWBD} changes to demand-model parameters.}
  \textcolor{red}{Panel (\textbf{a}) shows the multi-model mean change in SWBDs at 2.0$^\circ$C warming, with renewable share fixed at 50\%, the regional mix set to the current mix, and the \textcolor{red}{SWBD} threshold defined as the 99th percentile, under the default configuration: a cold comfort threshold of 12.5$^\circ$C, a hot comfort threshold of 19.6$^\circ$C, and heating/cooling response coefficients of 0.026/0.035. Panels (\textbf{b})--(\textbf{i}) show the [absolute/relative -- to confirm] difference from the default configuration, expressed as a percentage of reference average annual demand. Relative to the default configuration, panels \textbf{b} and \textbf{c} decrease and increase the cold threshold by 2$^\circ$C, respectively; panels \textbf{d} and \textbf{e} decrease and increase the hot threshold by 2$^\circ$C; panels \textbf{f} and \textbf{g} decrease and increase the response coefficients by 20\%; and panels \textbf{h} and \textbf{i} correspond to a wide and a narrow comfort zone, respectively. Black regions are excluded by the trend-validation filter; hatched regions indicate areas without current renewable capacity.}}
  \label{edf:demand_sensitivity}
\end{figure}

\vspace*{\fill}


\vspace*{\fill}

\end{document}